The principal comparison starts with FP32 queries in host memory and ends
when top-10 IDs and scores are available in host memory after stream
completion (\cref{fig:3-boundary}). It includes pinned staging, H2D,
query preprocessing, search and D2H. CAGRA int8 query quantisation is inside
this boundary; quantiser training and index construction are outside it.
Resource handles and output buffers are reused. JHQ returns to pageable
host arrays and the baselines to pinned arrays; CAGRA/PQ use Python APIs,
whereas JHQ/RaBitQ use C++. These interface differences remain visible when
interpreting small margins.

\par\medskip\noindent\begin{minipage}{\linewidth}
\centering\captionsetup{hypcap=false}
\resizebox{\linewidth}{!}{% Protocol diagram, not a measured stage-time plot.
\begin{tikzpicture}[
 font=\small\linespread{1}\selectfont,
 box/.style={draw=ThesisNavy!65,rounded corners=2pt,align=center,
 text width=2.55cm,minimum height=1.55cm,inner sep=4pt},
 arr/.style={-{Latex[length=2mm]},draw=ThesisNavy!80,thick}]
\node[box] (host) at (0,0) {Host FP32 queries\\Pinned staging};
\node[box,fill=ThesisNavy!7] (upload) at (3.5,0) {H2D transfer\\Query preprocessing};
\node[box,fill=ThesisNavy!7] (gpu) at (7,0) {GPU search\\Selection / refinement};
\node[box] (result) at (10.5,0) {D2H IDs and scores\\Stream completion};
\draw[arr] (host) -- (upload);
\draw[arr] (upload) -- (gpu);
\draw[arr] (gpu) -- (result);
\draw[ThesisNavy,thick] (-1.4,1.15) -- (-1.4,1.45) -- (11.9,1.45) -- (11.9,1.15);
\node[fill=white,inner sep=3pt] at (5.25,1.45) {Primary host-to-host timed interval};
\node[box,text width=5.7cm,minimum height=1.3cm] (setup) at (1.65,-2.35)
 {Outside search timing\\Load data / train or load model / add / warm-up};
\node[box,text width=5.7cm,minimum height=1.3cm] (cal) at (8.65,-2.35)
 {Separate policy cost\\Calibration; reuse cost model};
\draw[densely dashed,ThesisNavy!60] (-1.4,-1.1) -- (11.9,-1.1);
\end{tikzpicture}
}
\captionof{figure}[Host and device timing boundaries]{Timing boundaries. The host interval includes transfers and
query-specific computation. Device-stage profiles explain parts of that
interval; separate profiling and synchronisation can change execution, so
their sum is not an interchangeable end-to-end latency.}
\label{fig:3-boundary}
\end{minipage}\par\medskip

For round $r$ with $m_r$ timed calls, each answering a batch of $b$ queries,
\begin{equation}
Q_r=\frac{m_rb}{\sum_{i=1}^{m_r}t_{ri}},\qquad
Q=\operatorname{median}_r Q_r,
\label{eq:3-throughput}
\end{equation}
with $t$ in seconds. E1 uses $b=1{,}000$ and five rounds. JHQ performs
one untimed call per variant per round, then normally 24 timed calls.
Baselines perform three initial warm-up calls and normally 24 calls per
round, reducing that count for points classified as slow. Actual counts
are retained. Warm-up excludes JIT and graph capture from steady-state
timing. A batch's median or p95 latency describes repeated batch calls,
not the latency distribution of individual queries or concurrent requests.

The median of round throughputs is distinct from the reciprocal of median
call latency. Both summaries are retained where available, together with
round minima/maxima and baseline call traces. Slow valid calls remain in
the primary estimate. Rounds share a trained index, and budget draws share
a query pool: their counts measure timing variation or sample sensitivity,
rather than independent training replications. Small paired differences
are reported with their spread and consistency; no independence-based
significance claim is attached to the sweep cells.

The principal environment is an RTX 5090 with 32,607\,MiB reported memory,
driver 595.71.05, CUDA toolkit 13.0 and cuVS 26.08.01. Runs acquire a
cooperative GPU lock and record environment snapshots; OpenMP uses 16
threads within a 25-CPU container quota. E1 records a 575\,W power limit.
These snapshots do not establish locked clocks or constant host load.
JHQ's Release compute\_80 PTX is JIT-compiled \cite{cuda_platform}; library baselines use their
installed builds. Auxiliary device-only timing and stage profiles are
labelled separately, since JHQ's main API accepts host queries.

Memory accounting separates persistent encoded arrays, model/IVF metadata,
query workspace and runtime/context allocations. A process used-memory
sample includes more than the index; a build-time allocation delta excludes
memory pools created later. E1's context-subtracted JHQ totals and CAGRA's
index deltas retain these different definitions. Component models explain
resident totals and expose an unattributed remainder
(\cref{sec:4-4,sec:7-7}); they do not supply an equal-memory performance comparison.
Supplementary 20\,ms device samples separately measure construction peaks. Cold train/add timers likewise
exclude input loading and preceding object/host-model setup.

\paragraph{Supplementary evidence boundary.}
The supplementary campaigns add the checks catalogued in
\cref{sec:a-supplement}. Their hardware class and reported driver/toolkit
match the archives, but timing is a separate session. Arithmetic modes
and the cached-model precision intervention are specified in
\cref{sec:4-2,sec:7-precision}; historical baseline comparisons retain
that session distinction.
